\chapter{Regole di corrispondenza con ERA5}
\label{cap:regole}

Due regole decidono che cosa l'archivio pubblica e dove, e valgono per tutte le 90 variabili.

\section{Regola di pubblicazione}
\label{sec:pubblicazione}

Un campo si pubblica sotto un nome ERA5 solo se la grandezza WRF è quella che il parametro ERA5
definisce, oppure se la differenza si può dire in una riga. Un campo che non soddisfa nessuna
delle due condizioni resta fuori dall'archivio, anche quando il modello lo calcola. I campi esclusi
restano fuori per uno di due motivi, e non sono la stessa cosa: per alcuni campi WRF nessun
parametro ECMWF descrive la stessa grandezza, per altri il parametro esiste ma ERA5 non lo pubblica
(\cref{sec:non-pubblicati}).

Sette campi, in sei righe della \cref{tab:approssimazioni}, sono pubblicati con una differenza
dichiarata rispetto al loro omonimo ERA5.

\begin{table}[htbp]
  \centering
  \small
  \caption{Campi pubblicati con una differenza dichiarata rispetto al parametro ERA5.}
  \label{tab:approssimazioni}
  \begin{tabularx}{\textwidth}{@{}lXXl@{}}
    \toprule
    Campo & Parametro ERA5 & Campo CHAPTER & Sezione \\
    \midrule
    \campo{slor} & pendenza dell'orografia sottogriglia & pendenza del terreno risolto a \SI{3}{km} & \ref{sec:orografia} \\
    \campo{10fg} & raffica parametrizzata & massimo orario del vento risolto a \SI{10}{m} & \ref{cap:raffica} \\
    \campo{iews}, \campo{inss} & componenti dello sforzo turbolento in superficie dell'IFS & $\rho u_*^2$, con $\rho$ densità dell'aria e $u_*$ velocità d'attrito, proiettato sulla direzione del vento a \SI{10}{m} & \ref{sec:rugosita} \\
    \campo{fsr} & rugosità efficace, con la resistenza orografica parametrizzata & rugosità locale \campo{ZNT}, valore di tabella della categoria dominante modificato dalla neve & \ref{sec:rugosita} \\
    \campo{lsm} & maschera terra--mare permanente & maschera oraria di WRF, che conta come terra il ghiaccio marino & \ref{sec:statici-orari} \\
    \campo{tcw} & acqua totale in colonna, cinque specie & somma di sei specie, graupel incluso & \ref{sec:colonne} \\
    \bottomrule
  \end{tabularx}
\end{table}

\section{Regola di mascheratura}
\label{sec:mascheratura}

Un punto è mancante dove il parametro ERA5 non è definito. Un punto in cui il modello non fornisce
un valore, o in cui il valore sembra sbagliato, non per questo è mancante. I punti mancanti sono
scritti come bitmap GRIB, mai come zeri, e chi legge l'archivio distingue sempre un punto in cui il
parametro non è definito da uno in cui vale zero. La \cref{tab:maschere} elenca i campi mascherati.

\begin{table}[htbp]
  \centering
  \small
  \caption{Campi con punti mancanti e dominio su cui sono definiti.}
  \label{tab:maschere}
  \begin{tabularx}{\textwidth}{@{}lXl@{}}
    \toprule
    Campo & Definito su & Sezione \\
    \midrule
    \campo{sst}, \campo{ci} & celle di mare & \ref{sec:mare} \\
    \campo{dl} & celle con copertura di lago maggiore di zero (\num{32000}) & \ref{sec:vegetazione} \\
    \campo{swvl1-4}, \campo{stl1-4}, \campo{slt} & terra permanente (\num{1304109} celle) & \ref{sec:suolo}, \ref{sec:vegetazione} \\
    \campo{tvl} & tutte le celle, tranne le \num{124672} in cui domina un arbusteto & \ref{sec:vegetazione} \\
    \campo{mucin} & colonne con CAPE di almeno \SI{100}{J/kg} & \ref{sec:instabilita} \\
    \campo{tsn} & celle con \campo{SNOWC} sopra \num{0.9} & \ref{sec:neve} \\
    \campo{fal} & punti con radiazione solare incidente in superficie sopra \SI{50}{W/m^2} & \ref{sec:albedo} \\
    \bottomrule
  \end{tabularx}
\end{table}

La copertura dei laghi \campo{cl} è definita anche sul mare, dove vale zero, perché una cella di
mare non ha copertura lacustre. Quattro campi sono mascherati per una ragione diversa: nei punti mancanti
la grandezza WRF non ha un valore che corrisponda al parametro ERA5. \campo{tvl} manca sugli
arbusteti, perché la tabella dei tipi ERA5 separa arbusti sempreverdi e decidui e la classificazione
MODIS non porta la fenologia degli arbusti (\cref{sec:vegetazione}). \campo{tsn} manca dove la neve
non copre la cella, perché lì il campo del modello è una temperatura del suolo (\cref{sec:neve}).
\campo{fal} manca dove la radiazione solare incidente in superficie non supera \SI{50}{W/m^2};
senza radiazione incidente il rapporto fra radiazione riflessa e incidente non è definito (\cref{sec:albedo}). \campo{mucin} manca dove CAPE è sotto \SI{100}{J/kg}, perché lì la routine non restituisce
un'inibizione (\cref{sec:instabilita}).

\subsection{Mare e acque interne}
\label{sec:mare}

Per mascherare \campo{sst} e \campo{ci} serve distinguere il mare dalle acque interne, e la classe
di lago del file statico non basta. La classe 21 (lago) della classificazione MODIS-IGBP comprende
il Mar Nero, il Mare d'Azov e parte del Baltico, che insieme sono l'\SI{89}{\percent} delle
\num{55881} celle che il file statico classifica interamente come lago. Mascherare su quella classe
cancellerebbe \campo{sst} dal Mar Nero.

Il criterio usato è la dimensione della componente connessa, cioè dell'insieme di celle d'acqua
contigue per lato. Nel dominio le componenti più grandi hanno \num{844431} celle (Atlantico,
Mediterraneo, Baltico e Mare del Nord), \num{49841} (Mar Nero e Mare d'Azov) e \num{8583} (Mar
Rosso). La successiva è il lago Vänern, con \num{1225} celle, il lago più grande del dominio. Sono
mare le componenti di almeno \num{2000} celle, per un totale di \num{902855} celle di mare.
Qualunque soglia fra \num{1226} e \num{8583} celle dà lo stesso insieme.

\section{Campi statici e campi orari}
\label{sec:statici-orari}

Un campo statico viene dal file statico ed è quindi identico in ogni file orario. Sono statici
\campo{cl}, \campo{cvh}, \campo{cvl}, \campo{dl}, \campo{slt}, \campo{tvh} e \campo{tvl}. Le
colonne di suolo \campo{swvl1-4} e \campo{stl1-4} cambiano ogni ora, ma la loro maschera è la terra
permanente del file statico.

La maschera terra--mare \campo{lsm} è invece la maschera oraria di WRF (\campo{LANDMASK}), che
riclassifica come terra ogni cella d'acqua coperta da ghiaccio marino. La maschera cresce quindi in
inverno: in un file del 20 marzo 2024 le celle riclassificate sono \num{2469}, esattamente le celle
con ghiaccio marino. Il ghiaccio marino \campo{ci} vale 0 o 1, perché la simulazione non usa
ghiaccio frazionario (\campo{FRACTIONAL_SEAICE=0}). La maschera di terra permanente, nel senso di
ERA5, è l'insieme dei punti con $\campo{lsm} = 1$ e $\campo{ci} = 0$, contando come zero i punti
mancanti di \campo{ci}.

Anche l'albedo di fondo \campo{al} cambia nel tempo, con il ghiaccio marino e al passaggio fra la
stagione estiva e quella invernale della tabella di uso del suolo (\cref{sec:albedo}). Per
\campo{lsm} e \campo{al} vale il file orario dell'istante di validità, non un file qualunque
dell'archivio.
